% !TeX root = main.tex
\section{The derived category of a gerbe}
\label{sec:decomp-deriv-categ}
In this section we generalize some of the results by C{\u{a}}ld{\u{a}}raru
\cite[Chapter~2]{caldararu2000} and Lieblich \cite[Section~2.2.4]{lieblich2004}
on the basic structure of derived categories of gerbes and twisted sheaves.
Throughout the section,
we let $\alpha\colon X \to S$ be a gerbe banded by a diagonalizable group $\Delta$
with character group $A$.

\begin{definition}
\label{def:derived-twisted}
Given a character $\chi \colon \Delta_S \to \GGms{S}$,
we say that an object $\mathcal{F}$ in $\D_\qc(X)$ is
\define{$\chi$-homogeneous} if it has $\chi$-homogeneous
cohomology.
We denote the full subcategory of $\chi$-homogeneous objects in $\D_\qc(X)$ by~$\D_{\qc, \chi}(X)$.

Consider the special case $\Delta = \GGm$ and let $\id\colon \GGm \to \GGm$ be the identity character.
An object in $\D_{\qc, \id}(X)$ is called an \emph{$\alpha$-twisted complex on $S$}.
We often use the notation $\D_{\qc, \alpha}(S)$ instead of~$\D_{\qc, \id}(X)$
employing a similar abuse of notation as in Remark~\ref{rem-twisted-abuse}
if $\alpha$ is just a cohomology class in $\HH^2(S_\fppf, \GGm)$.
\end{definition}

We state the following lemma for convenient reference.
See Remark~\ref{rem-trivial-case} for sharper results.

\begin{lemma}
\label{lem:affine-derived-decomp}
Let $S$ be an affine scheme and let $\alpha\colon X \to S$ be a trivial $\Delta$-gerbe,
where  $\Delta$ is a diagonalizable group.
Then the obvious functor
\begin{equation}
\label{eq-equiv-qcoh}
\D(\Qcoh(X)) \xrightarrow{\sim} \D_\qc(X)
\end{equation}
is a monoidal equivalence and induces equivalences
\begin{equation}
\label{eq-equiv-qcoh-chi}
\D(\Qcoh_\chi(X)) \xrightarrow{\sim} \D_{\qc, \chi}(X)
\end{equation}
for each character $\chi$ of $\Delta$.
Moreover, let $T$ be an affine scheme and let $\beta\colon X' \to T$ be a trivial $\Delta'$-gerbe,
where $\Delta'$ is a diagonalizable group.
Assume that we have a morphism $\rho\colon X \to X'$ of
algebraic stacks.
Then the squares
\begin{equation}
\label{eq-commute-derived}
\xymatrix{
\D(\Qcoh(X')) \ar[r]^-{\sim}\ar[d]_{\Ld\rho^*} & \D_\qc(X') \ar[d]^{\rho^*}\\
\D(\Qcoh(X)) \ar[r]^-{\sim}& \D_\qc(X), \\
}\qquad
\xymatrix{
\D(\Qcoh(X')) \ar[r]^-{\sim} & \D_\qc(X') \\
\D(\Qcoh(X)) \ar[r]^-{\sim}\ar[u]^{\Rd\rho_*}& \D_\qc(X)\ar[u]_{\rho_*} \\
}
\end{equation}
commute up to natural equivalence.
\end{lemma}

\begin{proof}
The equivalence \eqref{eq-equiv-qcoh} is a special case of \cite[Proposition~2.2.4.6]{lieblich2004},
and the equivalences \eqref{eq-equiv-qcoh-chi} are then immediate
from Theorem~\ref{t-qc-gerbe-split}.
The fact that the equivalence \eqref{eq-equiv-qcoh} is monoidal follows from the fact that
any complex in $\Qcoh(X)$ admits a resolution 
by an h-flat complex in $\Qcoh(X)$ which is also h-flat as a
complex in $\Mod(X_\liset, \mathcal{O}_X)$.
This can be seen by using a homotopy limit argument (see~\cite[\sptag{06XX}]{stacks-project}),
using that any object in $\Qcoh(X)$ is a quotient of a locally free sheaf and the fact that coproducts in $\Mod(X_\liset, \mathcal{O}_X)$ are exact and preserve $\Qcoh(X)$.

The functor $\alpha_*\colon \Qcoh(X) \to \Qcoh(S)$ is exact,
since this operation corresponds to taking invariants and $\Delta$ is diagonalizable.
This well-known fact also easily follows from the last sentence of Proposition~\ref{t:equivalence-epi}
together with the decomposition in~Theorem~\ref{t-qc-gerbe-split}.
Hence the equivalence \eqref{eq-equiv-qcoh} shows that $X$ has cohomological dimension zero.

Note that $\rho$ induces a morphism $S \to T$ on the coarse spaces,
so $\rho$ factors as $X \to X'\times_T S \to X'$ where the first morphism
is a morphism of gerbes over $S$, and the second morphism is affine.
Furthermore, the first morphism factors as an affine morphism followed by a
$\Delta$-gerbe by Remark~\ref{rem-gerbe-ff}.
Hence also $\rho$ has cohomological dimension zero,
as this property is stable under composition and can be checked after a faithfully flat base change
by \cite[Lemma~2.2]{hr2017}).
In particular, $\rho$ is concentrated, so it follows from \cite[Corollary~2.2]{hnr2018} that the right hand
square \eqref{eq-commute-derived} commutes.
Hence also the left hand square commutes by adjunction.
\end{proof}

\begin{remark}
\label{rem-gerbe-concentrated}
Let $\rho\colon X \to Y$ be a morphism of gerbes banded by diagonalizable groups
over an algebraic stack $S$ as in Remark~\ref{rem-gerbe-ff}.
Then the last part of the proof shows that $\rho$ is concentrated of cohomological dimension zero,
as these properties can be checked locally on~$S$
(see \cite[Lemma~2.2 and Lemma~2.5]{hr2017}).
\end{remark}

The next theorem is the main structure theorem for the derived
category of a $\Delta$-gerbe and the derived analog of Theorem~\ref{t-qc-gerbe-split}.
It is a generalization of the observation in \cite[Section~2.2.4]{lieblich2004}
that the derived category splits according to characters in certain situations.
\begin{theorem}
  \label{t:chi-splitting-D-qc-X}
  Let $S$ be an algebraic stack and 
  $\alpha\colon X \to S$ a gerbe banded by a diagonalizable group
  $\Delta$.
  Then $\D_{\qc, \chi}(X)$ is
  a triangulated subcategory of $\D_{\qc}(X)$ for each $\Delta_S$-character $\chi$,
  and taking the coproduct gives an equivalence
\begin{equation}
\label{eq-equivalence-derived}
\prod_{\chi \in A} \D_{\qc, \chi}(X) \sira \D_\qc(X),
\qquad
 (\mathcal{F}_\chi)_{\chi \in A} 
\mapsto 
\bigoplus_{\chi \in A} \mathcal{F}_\chi,
\end{equation}
of triangulated categories, where $A$ denotes the character group of $\Delta$.

Furthermore, assume that $\mathcal{F}_\chi$ and $\mathcal{G}_\psi$ are objects in $\D_\qc(X)$
which are homogeneous for $\Delta_S$-characters $\chi$ and $\psi$, respectively.
Then the objects
\begin{equation}
\label{eq-der-monoidal}
\mathcal{F}_\chi \otimes \mathcal{G}_\psi, \qquad
\sheafHom(\mathcal{G}_\psi, \mathcal{F}_\chi)
\end{equation}
are $\chi\psi$- and $\chi\psi^{-1}$-homogeneous, respectively.
\end{theorem}

\begin{proof}
By Theorem \ref{t-qc-gerbe-split},
the category $\Qcoh_\chi(X)$ is a Serre subcategory of $\Qcoh(X)$.
It follows that $\D_{\qc, \chi}(X)$ is a strictly full triangulated subcategory of $\D_\qc(X)$
by \cite[\sptag{07B4}, \sptag{06UQ}]{stacks-project}.

Recall that we work with sheaves on the lisse--étale sites over $X$ and~$S$.
Note that we get a morphism $X_{\liset} \to S_{\liset}$ of topoi,
by 
\cite[\sptag{07AT}]{stacks-project},
since $\alpha$ is smooth (see Remark~\ref{rem-gerbe-ff}).
In particular, there is no need to work with hypercoverings to define pull-backs as in~\cite{olsson2007}.

Let $E$ be a subset of the set $A$ of characters of $\Delta$,
and let $\Qcoh_E(X)$ denote the full subcategory of $\Qcoh(X)$ of objects
which are direct sums of $\chi$-homogeneous objects for $\chi \in E$.
Let $F$ be the complement of $E$ in $A$.
It follows from Theorem~\ref{t-qc-gerbe-split}
that $(\Qcoh_E(X), \Qcoh_F(X))$ is a torsion pair in $\Qcoh(X)$ and that both
$\Qcoh_E(X)$ and~$\Qcoh_F(X)$ are abelian subcategories of~$\Qcoh(X)$.

Let $\mathcal{E}$ and $\mathcal{F}$ be objects of $\Qcoh_E(X)$ and~$\Qcoh_F(X)$,
respectively.
We will show that
\begin{equation}
\label{eq-ext-vanishing}
\Ext^n_{\mathcal{O}_X}(\mathcal{E}, \mathcal{F})
\cong \HH^n(\Rd\Gamma\Rd\sheafHom_{\mathcal{O}_X}(\mathcal{E}, \mathcal{F}))
\end{equation}
vanishes for each $n$.
Since
$
\Rd\Gamma\Rd\sheafHom_{\mathcal{O}_X}(\mathcal{E}, \mathcal{F})
\cong
\Rd\Gamma\Rd\alpha_\ast\Rd\sheafHom_{\mathcal{O}_X}(\mathcal{E}, \mathcal{F}),
$
this follows if we prove that $\Rd\alpha_\ast\Rd\sheafHom_{\mathcal{O}_X}(\mathcal{E}, \mathcal{F})$ vanishes.
The $n$-th cohomology sheaf of $\Rd\alpha_\ast\Rd\sheafHom_{\mathcal{O}_X}(\mathcal{E}, \mathcal{F})$ is easily seen
to be the sheaf of $\mathcal{O}_S$-modules associated to the presheaf
\begin{equation*}
(U \to S) \mapsto
\Ext^n_{\mathcal{O}_{X_U}}(\mathcal{E}|_{X_U}, 
\mathcal{F}|_{X_U})
\end{equation*}
where $U$ is a scheme and $U \ra S$ is smooth. 
It is therefore enough to verify the vanishing of \eqref{eq-ext-vanishing} when $S$ is affine and the gerbe
$\alpha\colon X \to S$ is trivial.
But in this situation the functor \eqref{eq-equiv-qcoh} is an equivalence by Lemma~\ref{lem:affine-derived-decomp},
so we get the desired vanishing from
Lemma~\ref{p:abelian-torsion-pair}.\ref{enum:Ext-TF} or, more
easily, from
Theorem~\ref{t-qc-gerbe-split}.

All direct and inverse truncation systems in $\D_\qc(X)$ are
effective by Propositions~\ref{p:ab4-trunc} and \ref{p:Dqc-truncation-complete}.
Hence the hypotheses of Theorem~\ref{t:sod-DBA} are satisfied,
so we get a semi-orthogonal decomposition
\begin{equation}
\label{eq-so-decomp}
\D_\qc(X) = \langle \D_{\qc, F}(X), \D_{\qc, E}(X)\rangle,
\end{equation}
where $\D_{\qc, E}$ denotes the full subcategory of $\D_\qc(X)$ of objects with cohomology in
$\Qcoh_E(X)$ and similarly for $\D_{\qc, F}(X)$.

We are now ready to prove that the functor \eqref{eq-equivalence-derived} is an equivalence.
First note that it is well defined since $\D_\qc(X)$ has arbitrary coproducts.
Let $\chi \in A$ and set
$E = \{\chi\}$, $F = A -\{\chi\}$.
Then $\Hom_{\D_\qc(X)}(\mathcal{F}_\chi, \mathcal{G})$ vanishes
whenever $\mathcal{F}_\chi \in \D_{\qc, \chi}(X)$
and $\mathcal{G} \in \D_{\qc, F}$ by the 
semi-orthogonal decomposition~\eqref{eq-so-decomp}.
Since $\D_{\qc, \chi}(X)$ is a full subcategory of $\D_\qc(X)$,
it follows that the functor \eqref{eq-equivalence-derived} is fully faithful.

Now let $\mathcal{F}$ be any object of $\D_\qc(X)$.
Then the semi-orthogonal decomposition~\eqref{eq-so-decomp}
gives a triangle
$$
\mathcal{F}_\chi \to \mathcal{F} \to \mathcal{F}_F \to \Sigma\mathcal{F}_\chi
$$
with $\mathcal{F}_\chi \in \D_{\qc, \chi}(X)$ and $\mathcal{F}_F
\in \D_{\qc, F}(X)$, respectively.
By Remark~\ref{rem:cohomology-split} 
and Theorem~\ref{t-qc-gerbe-split},
the morphism $\mathcal{F}_\chi \to \mathcal{F}$ induces an isomorphism
$\HH^n(\mathcal{F}_\chi) \to \HH^n(\mathcal{F})_\chi$ for each $n$.
It follows that the canonical map $\bigoplus_{\chi \in A} \mathcal{F}_\chi \to \mathcal{F}$ induces
an isomorphism on cohomology and therefore is an isomorphism in $\D_\qc(X)$.
Hence the functor \eqref{eq-equivalence-derived} is essentially
surjective and therefore an equivalence. 

Finally, we prove the statements about the monoidal structure.
It is enough to prove the statement about the tensor product,
since the statement about the internal hom follows formally from the
adjointness property.
The question whether $\mathcal{F}_\chi \otimes_\mathcal{O_X}^\Ld \mathcal{G}_\psi$
is $\chi\psi$-homogeneous can be verified locally on $S$.
Indeed, this follows from Proposition~\ref{p:pull-back-grading},
since cohomology of a complex commutes with pull-back along a flat map (see \cite[Equation~1.9]{hr2017}),
and the pull-back functor along a faithfully flat map is conservative.
Hence we may assume that $S$ is affine and that the gerbe $\alpha\colon X \to S$
is trivial.
Now the statement follows from the underived case described in Theorem~\ref{t-qc-gerbe-split}
by the monoidal decomposition preserving
equivalence \eqref{eq-equiv-qcoh} in Lemma~\ref{lem:affine-derived-decomp}.
\end{proof}

\begin{remark}
Similarly as in Remark~\ref{rem-coprod-prod}, the canonical morphism
$
\bigoplus_{\chi\in A} \mathcal{F}_\chi
\to
\prod_{\chi\in A} \mathcal{F}_\chi
$
is an isomorphism even if the character group
$A$ is not finite.
\end{remark}

The following two propositions are the derived analogs of
Proposition~\ref{p:pull-back-grading} and \ref{t:equivalence-epi}.

\begin{proposition}
\label{p:pull-back-grading-derived}
Let $S$ be an algebraic stack and 
$\alpha\colon X \to S$ a gerbe banded by a diagonalizable group
$\Delta$.
Let $f\colon T \to S$ be a morphism of algebraic stacks,
and let $g\colon Y \to X$ denote the canonical morphism from the pull-back
$Y$ of $X$ along $f$.
Then the functors
$$
g^*\colon \D_\qc(X) \to \D_\qc(Y), \qquad
g_*\colon \D_\qc(Y) \to \D_\qc(X)
$$
respect the decompositions of $\D_\qc(X)$ and~$\D_\qc(Y)$
into homogeneous components given in Theorem~\ref{t:chi-splitting-D-qc-X}.
\end{proposition}
\begin{proof}
As a formal consequence of adjunction,
it is enough to verify that the pull-back $g^*$ respects the decomposition.
By a similar reduction argument as in the last paragraph of
the proof of Theorem~\ref{t:chi-splitting-D-qc-X},
we reduce to the case when $S$ and $T$ are affine and the gerbes are trivial.
Now the result follows from the underived setting in Proposition~\ref{p:pull-back-grading}
and the commutative diagram \eqref{eq-commute-derived} in Lemma~\ref{lem:affine-derived-decomp}.
\end{proof}

\begin{proposition}
  \label{t:equivalence-epi-derived}
  Let $X$ and $X'$ be gerbes over $S$ banded by diagonalizable groups $\Delta$ and~$\Delta'$,
  respectively, and let $\rho\colon X \to X'$ be a morphism over~$S$ with induced morphism $\varphi\colon \Delta_S \to \Delta'_S$ on the bands.
  Fix a character $\chi'\colon \Delta'_S \to \GGms{S}$, and let $\chi =
  \chi' \circ \varphi$.
  Then the pull-back functor $\rho^*$ takes $\chi'$-homogeneous
  objects to $\chi$-homogeneous 
  objects and induces an equivalence
  \begin{equation}
    \label{eq:restricted-pb-derived}
    \rho^*\colon \D_{\qc,\chi'}(X') \sira \D_{\qc,\chi}(X)
  \end{equation}
  of triangulated categories.
  In particular, if $X' = S$ then $\rho^*\colon \D_\qc(S) \to \D_\qc(X)$ is fully faithful with 
  essential image consisting of those objects that are
  homogeneous for the trivial character. 
\end{proposition}

\begin{proof}
Since $\rho$ is flat (see Remark~\ref{rem-gerbe-ff}),
the pull-back $\rho^*$ commutes with
taking cohomology (see~\cite[Equation~(1.9)]{hr2017}).
In particular,
the first statement follows from the underived case,
which is given by Proposition~\ref{t:equivalence-epi}.

Let us first consider the special case $X' = S$.
By Remark~\ref{rem-gerbe-concentrated}, the morphism $\rho$ is concentrated,
so the push-forward $\rho_*$ respects flat base change by \cite[Theorem~2.6(4)]{hr2017}.
Hence the statement can be verified locally.
In particular, we may assume that the gerbes are trivial and that $S$ is affine.
Hence we may use the equivalences in Lemma~\ref{lem:affine-derived-decomp}.
Since $\rho^*$ and $\rho_*$ are exact,
the statement follows from the underived case,
which is Proposition~\ref{t:equivalence-epi}.

For the general case, we proceed exactly as in
the proof of Proposition~\ref{t:equivalence-epi} on the level of derived
categories.
\end{proof}

Similarly as in the underived case, we get the following corollaries.
\begin{corollary}
\label{cor:twisted-power-derived}
Let $S$ be an algebraic stacks and let $\alpha\colon X \to S$ be a gerbe banded by $\Delta$,
and choose a character $\chi\colon \Delta_S \to \GGms{S}$.
Then we have an equivalence
\begin{equation*}
%\label{eq-pb-character-derived}
\D_{\qc,\chi_*\alpha}(S)\cong \D_{\qc, \chi}(X) 
\end{equation*}
of categories, where $\chi_*\alpha$ is defined as in Construction~\ref{cons-gerbe-pf}.
In the particular case when $\Delta = \GGm$,
this gives the equivalence
\begin{equation*}
%\label{eq-pb-power-derived}
\D_{\qc,\alpha^d}(S)\cong \D_{\qc,d}(X) 
\end{equation*}
for any $d \in \mathbb{Z}$, where $\alpha^d$ is defined as in Construction~\ref{cons-gerbe-operations}.
\end{corollary}
\begin{proof}
This is proven similarly as Corollary~\ref{cor:twisted-power} by using Proposition~\ref{t:equivalence-epi-derived}
instead of Proposition~\ref{t:equivalence-epi}.
\end{proof}

\begin{corollary}[{cf.~\cite[Theorem 2.2.4]{caldararu2000}}]
\label{cor:twisted-tensor-derived}
Let $S$ be an algebraic stack and let $\alpha\colon X \to S$ and~$\beta\colon Y \to S$ be $\GGm$-gerbes.
Then the tensor product and the internal Hom-functor on $\D_\qc(X\times_S Y)$ induce triangulated functors
\begin{align*}
\otimes & \colon \D_{\qc, \alpha}(S) \times \D_{\qc, \beta}(S)
\to \D_{\qc, \alpha\beta}(S),\\
\sheafHom & \colon \D_{\qc, \alpha}(S)^\opp \times \D_{\qc, \beta}(S)
\to \D_{\qc, \alpha^{-1}\beta}(S).
\end{align*}
Furthermore, if $f\colon T \to S$ is an arbitrary morphism of algebraic stacks,
then we get an induced pair of adjoint functors
\begin{equation*}
  f^* \colon \D_{\qc, \alpha}(S) \to \D_{\qc, f^*\alpha}(T),
  \qquad
  f_* \colon \D_{\qc, f^*\alpha}(T) \to \D_{\qc, \alpha}(S).
\end{equation*}
\end{corollary}
\begin{proof}
This is proven similarly as Corollary~\ref{cor:twisted-tensor} by using 
Theorem~\ref{t:chi-splitting-D-qc-X},
Proposition~\ref{p:pull-back-grading-derived}
and~Proposition~\ref{t:equivalence-epi-derived}
instead of
Theorem~\ref{t-qc-gerbe-split},
Proposition~\ref{p:pull-back-grading}
and~Proposition~\ref{t:equivalence-epi}.
\end{proof}

We conclude the section with two remarks regarding alternative
approaches for proving Theorem~\ref{t:chi-splitting-D-qc-X}.
\begin{remark}
\label{rem-trivial-case}
Note that Theorem~\ref{t:chi-splitting-D-qc-X} is trivial whenever the functor \eqref{eq-equiv-qcoh}
is a monoidal equivalence.
By Lieblich \cite[Proposition~2.2.4.6]{lieblich2004},
the functor \eqref{eq-equiv-qcoh} is an equivalence whenever $S$ is a quasi-compact,
separated scheme.
More generally, it follows from the work by Hall, Neeman and~Rydh \cite[Theorem~1.2]{hnr2018}
that \eqref{eq-equiv-qcoh} is an equivalence whenever
$S$ is an algebraic space which is either noetherian or quasi-compact with affine diagonal.
The aforementioned theorem applies since in this case,
the category $\D_\qc(X)$ is compactly generated by $|A|$ objects by
\cite[Theorem~6.9]{hr2017} 
applied to the presheaf
$
(T \to S) \mapsto \D_\qc(X_T)
$
of triangulated categories.
Since we are not going to use this result,
we leave the details of verifying the hypotheses of the
cited theorem to
the reader (cf.~\cite[Example~9.2]{hr2017}).
We do not know whether the equivalence \eqref{eq-equiv-qcoh} is monoidal in this generality.

Finally, it should be emphasized that there certainly are interesting cases when
\eqref{eq-equiv-qcoh} is not an equivalence. 
For instance, consider the $\GGm$-gerbe in Example~\ref{ex-gerbes}.\ref{it-universal-projective}.
The functor \eqref{eq-equiv-qcoh} is not an
equivalence for $X = \BB\GL_n$ if $n \geq 2$
by \cite[Theorem~1.3]{hnr2018},
so Theorem~\ref{t:chi-splitting-D-qc-X} seems to be non-trivial even in this basic case.
\end{remark}

\begin{remark}
\label{rem:decomposition-O-modules}
Theorem~\ref{t:chi-splitting-D-qc-X} would be trivial if
the category $\Mod(X_\liset, \mathcal{O}_X)$ admitted a decomposition
similar 
to the one for $\Qcoh(X)$ described in Theorem~\ref{t-qc-gerbe-split}.
We do not know whether this is the case,
but we suspect it is not.

In \cite[Exposé 1, Proposition~4.7.2]{sga-3-1-2011}, it is stated
that the category of (not necessarily quasi-coherent) sheaves of $G-\mathcal{O}_S$-modules
on the small Zariski site
is equivalent to the category of sheaves of $\mathcal{A}$-comodules for any affine group scheme $G = \Spec_S \mathcal{A}$ over
a scheme $S$.
If this were true, then it should be possible to adapt the proof
to our situation and to get a decomposition of $\Mod(X_\liset, \mathcal{O}_X)$.
However, we suspect that the previously cited proposition might be wrong, 
since it uses \cite[Exposé 1, Proposition~4.6.4]{sga-3-1-2011}
which in turn uses the fact that the canonical morphism
\begin{equation}
  \label{eq:qcoh-push-pull}
  \mathcal{F}\otimes_S \mathcal{A} \to \alpha_\ast\alpha^\ast\mathcal{F}
\end{equation}
is an isomorphism whenever $\mathcal{F}$ is quasi-coherent and
$\alpha\colon X = \Spec_S \mathcal{A} \to S$ is affine.
But this is certainly not true in general if we drop the condition that $\mathcal{F}$ be quasi-coherent. 

Let for example $S$ be a scheme with an open subscheme $U$
such that the inclusion $j \colon U \ra S$ 
is affine and quasi-compact
and $\mathcal{F}=j_! \mathcal{O}_U$ is not
quasi-coherent. Then the canonical morphism $j_!
\mathcal{O}_U \cong \mathcal{F} \otimes j_*\mathcal{O}_U \ra
j_*j^* \mathcal{F} \cong j_*\mathcal{O}_U$ is not an
isomorphism because $j_*\mathcal{O}_U$ is quasi-coherent.
Let $G:=S \sqcup U$ be the disjoint union. Then
$\alpha \colon G \ra S$ is in the obvious way an
affine \'etale group
scheme over $S$. In this case 
$\mathcal{A}=\alpha_*\mathcal{O}_G=
\mathcal{O}_S \oplus
j_*\mathcal{O}_U$ and the morphism
\eqref{eq:qcoh-push-pull} is not an isomorphism for 
$\mathcal{F}=j_! \mathcal{O}_U$.
\end{remark}

%%% Local Variables: 
%%% mode: latex
%%% TeX-master: "main"
%%% End: 
